% ============================================================================
\clearpage
\section*{v4.10.0-r28 Source-Only: Status of the Residual Bend and of the Counting Rule after the Explorations of 16--17 November}
\addcontentsline{toc}{section}{v4.10.0-r28 Source-Only: Status of the Residual Bend and of the Counting Rule after the Explorations of 16--17 November}
\label{sec:v4100-r28-bend-counting-status}
% ============================================================================

\paragraph{Purpose.}
Between v4.10.0-r27 and this layer the author asked for a series of tests and explorations on two open problems: the residual bend of the
lepton phase (OP3) and the counting rule of the horizon readout (OP1).  They were run as pre-registered tests H106--H111 and as explorations
without an H number, and were recorded only in notes.  This layer collects their status in the Companion.  It adds no new computation; every
number below is taken from the cited script output.  Earlier layers are not rewritten.  The notes are listed in
\S\ref{subsec:v4100-r28-notes}.

\noindent\textbf{Source-only status.}
This v4.10.0-r28 layer is source-only.  It does not build a PDF, does not assign a new DOI,
and does not authorize a public release.  The v5.1.0 Version DOI
\href{https://doi.org/10.5281/zenodo.24237112}{10.5281/zenodo.24237112} and the Concept DOI
\href{https://doi.org/10.5281/zenodo.19210728}{10.5281/zenodo.19210728} are unchanged by it.

\subsection*{v4.10.0-r28 results}
\label{subsec:v4100-r28-results}

\paragraph{Remark v4.10.0-r28.1 (six pre-registered tests on the carrier side, all refuted).}
Each test was fixed in a pre-registration note before the computation; the target was the pole-mass phase \(3\delta=0.666674\) within
\(2\cdot10^{-4}\).
\begin{center}
\small
\begin{tabular}{@{}>{\raggedright\arraybackslash}p{0.12\textwidth}>{\raggedright\arraybackslash}p{0.50\textwidth}>{\raggedright\arraybackslash}p{0.30\textwidth}@{}}
\toprule
test & mechanism & outcome\\
\midrule
H106 & singlet admixture to the spin-1 triangle & refuted; nearest phase \(0.640\) or \(0.679\)\\
H107 & nested normalisation (ten mechanisms, two solutions each) & refuted; no hit within \(2\cdot10^{-4}\)\\
H108 & entangled generations (second Majorana point) & refuted; phases far from the target\\
H109 & generation-dependent frames under cyclic assignment & refuted; the assignment respects generation symmetry and does not bend\\
H110 & one algebra per frame, the fourth as observer (HS(7) half) & refuted; pure \(K=2/3,\ \varphi=\pi/4\), singlet \(K=1/2,\ \varphi=0\), no variant at the target\\
H111 & time octant through the frames & refuted; a design flaw (the cross-Gram matrix in \(\C^2\) has rank two) was recognised after the computation and is
stated in the result note\\
\bottomrule
\end{tabular}
\end{center}
Two further explorations belong here.  Running masses at \(M_Z\) give \(3\delta=0.663370\), \(0.0033\) below the pole-mass phase; the
pole-mass convention of the rounds is kept.  Multiplying the Dixon algebra by the sedenions gives
\(\C\otimes\Hh\otimes\Oo\otimes\mathbb S\) of real dimension \(2\cdot4\cdot8\cdot16=2^{10}\), that is the level HS(10), not HS(7): every factor
multiplies the dimension.

\paragraph{Remark v4.10.0-r28.2 (diagnosis: the phase is rigid in the carrier geometry; correction).}
Symmetric constructions on the carrier fix the Bargmann phase by their overlap: spin \(\tfrac12\), \(1\), \(\tfrac32\), \(2\) give
\(0.785\), \(0.679\), \(0.650\), \(0.637\).  The value \(\tfrac23\) is not a rational part of \(2\pi\).  A first diagnosis concluded that the
second number (\(\delta\)) is a property of the field state and needs a selection law, and that further carrier variants would yield nothing.
That diagnosis missed the guarded core and the open prediction H102 (v4.10.0-r14.5--r14.7) and is corrected here: the solid angle
\(2\varphi=1.33333\) is the octant \(\pi/2\) minus a cap of solid angle \(\gamma_M=0.237533\) (deviation \(6.7\cdot10^{-5}\) sr), so
\(\varphi=\pi/4-\gamma_M/2\) has a geometric reading.  The core escapes the rigidity because it changes the connection on the generation
sphere, not the states.  No native dephasing or entanglement parameter \(\kappa\) was found; in the core picture the bend starts at the
octant and the deficit to be explained is \(\pi/4-\tfrac23=0.1187\).  The chain then reads
\[
\varphi\ \leftarrow\ \text{core}\ \leftarrow\ \Omega=\gamma_M\ \leftarrow\ \text{counting rule of frame M}\ \leftarrow\ \text{OP1},
\]
and the selection law for the state reduces to one question: why frame M counts every spin with weight two.

\paragraph{Remark v4.10.0-r28.3 (what H102 rests on).}
H102 predicts \(m_\tau=1777.134\) MeV, \(+2.3\sigma\) above PDG 2024 (\(1776.93\pm0.09\) MeV, unchanged in PDG 2026); Belle II alone is
compatible (\(+0.3\sigma\)).  The measured inputs are correct.  The critical input is \(\gamma_M\): H102 and the frame factor
\(\gamma_L=1.8647\) stand on the U(1) counting of Domagala--Lewandowski and Meissner (weight 2 per spin, \(\gamma=0.237533\)).  With the
SU(2) counting (\(2j+1\), \(\gamma=0.274067\)), which the later literature prefers for horizons, H102 would give \(m_\tau\approx1868\) MeV and
fail.  The value also depends on the form of the area spectrum: with \(j\) or \(j+\tfrac12\) in place of \(\sqrt{j(j+1)}\) the U(1) count
gives \(0.3497\) or \(0.2206\).  H102 is therefore a joint test of the core reading, the U(1) counting and the spectrum \(\sqrt{j(j+1)}\).

\paragraph{Remark v4.10.0-r28.4 (the U(1) axis has a native form; (S1) as a lock between two SU(2)).}
\begin{enumerate}[leftmargin=2em]
\item[(a)] The duality rotation \(z\mapsto e^{ia}z\) turns \(\chi\) about its helicity axis by \(-2a\) (200 random fields, deviation below
\(10^{-15}\)); the tachyon map is its element \(a=\pi/2\).  Its poles are the self-dual and anti-self-dual fields, which are the compensator
and partner ports (r16.2(b), r16.4(b)).  Frame M is defined through this axis (convention F).  The U(1) of the counting thus has a native
form: electromagnetic duality, with the helicity axis as its axis.
\item[(b)] The spatial rotation group, whose Casimir gives the area of a puncture, leaves \(s_\pm\) fixed (500 rotations, deviation
\(1.5\cdot10^{-14}\)); \(\chi\) is a singlet of it.  The SU(2) in which the frame maps act as Pauli matrices is internal.  The identification
(A) of r24 therefore equates two different groups and needs a lock between them.
\item[(c)] A natural lock is (S1\('\)): an elementary strand is a massless quantum crossing the horizon along the normal.  Then
\(J_\nu=\lambda\), only \(\pm j\) are allowed, and the DL counting follows as the state count of massless crossings.  Under (S1\('\)),
however, the frame maps do not act as rotations at a fixed normal, so (A) fails in the form stated in r24.  In the literature the U(1) axis
of the ABK construction is the horizon normal, counting only horizon states gives the U(1) reduction (Mitra), and \(S=A/4\) holds for every
\(\gamma\) in the spin-foam derivation (Bianchi); none of this fixes the weight two.
\end{enumerate}

\paragraph{Remark v4.10.0-r28.5 (a time-dependent correction term and an impermeable boundary).}
On the author's idea a correction term that bends the phase along the time evolution was fitted, with two methods fixed in advance
(\texttt{Methoden\_\allowbreak Notiz.md}, \texttt{Methoden\_\allowbreak Notiz\_\allowbreak Verbindung.md}).  As a correction of the states it can
reveal no dynamics: in the representation on \(S\) (r16.8, r16.9) \(K\) and \(\varphi\) depend only on the populations, no native element
shifts them, and the minimal path is the tangent of the Koide circle.  As a time-dependent connection that removes the flux
\(\Delta\Phi=0.118724\) from the octant triangle with \(K=2/3\) exact, a local economical flow empties a seam along the edges on three grids
and does not produce the core of H102; that is an exclusion.  An impermeable boundary with one puncture would be a \emph{sufficient}
condition for a core of the H102 kind, with free position and the target value as its area.  The stock does not justify such a boundary,
and the mirror structure of the edges does not either (it is a pseudoscalar).  The line ends at the same two points: why exactly one core,
and why \(\Delta\Phi=\gamma_M/2\).

\paragraph{Remark v4.10.0-r28.6 (topology of the core; several reflectors; two rasters).}
\begin{enumerate}[leftmargin=2em]
\item[(a)] On \(S^2\) one core is the only rigid bend; with the reflector \(\Theta=i\sigma_yK\) (antipodal, \(\Theta^2=-1\)) as a symmetry
it is a core pair \(\pm z_c\), one puncture in \(\mathbb{RP}^2\).  A field with \(F(-n)=F(n)\) forces a core pair, and Borsuk--Ulam forces a
flux.  The size \(\Delta\Phi=\gamma_M/2\) is not topologically reachable, since \(\tfrac23\) is no rational part of \(2\pi\).
\item[(b)] Compared with the cryptographic sections, there is more than one reflector: \(C\) (\(z\mapsto-z\), the identity on \(\chi\)),
\(J_B\), and the rotor conjugates \(R_\alpha^{-1}J_BR_\alpha\); at \(\alpha=\pi/8\) core and antipode are exchanged.
\item[(c)] Two rasters were compared: the triole raster (\(D_3\), \(120^\circ\), generation cycle) and the quarter raster (\(D_4\)), which
contains the Klein group of the frame maps.  The axes of the effective reflectors lie at the magic angle.  Neither raster fixes
\(\gamma_M\); H111 was the test of this line (Remark r28.1).
\end{enumerate}

\paragraph{Remark v4.10.0-r28.7 (counting per frame: blurred readout, a separate HC per frame, register entropy).}
\begin{enumerate}[leftmargin=2em]
\item[(a)] A counting weight two per spin is compatible with full rotation invariance if the readout averages over the axis
(acceptance \(2/(2j+1)\)); the objection ``U(1) is a gauge fixing'' then does not apply.  The HC selects the states that saturate the
uncertainty relation (v4.9.0-r63.2).  A counter-HC read as complement leads to the SU(2) value \(0.274\); read as blurring it keeps \(0.2375\).
\item[(b)] With one HC per frame the counts nest as \(T\subseteq M\subseteq P\): T counts \(j=\tfrac12\) (\(0.127384\)), M weight two
(\(0.237533\)), P all \(2j+1\) (\(0.274067\)).  The last is a setting motivated by the degeneracy of the ports (F2).
\item[(c)] In convention F, P lies between M and T in the coordinate \(x\).  In rapidity, \(x=\tanh^2\eta\), M lies at \(\eta\) and T at
\(\eta+i\pi/2\) (\(1/x\)); both branches end at P for \(\eta\to\infty\).  P is the common endpoint, and seen from R it lies outside, as in the
counting.  A limit of rank-two masks never reaches \(2j+1\), so the HC of P is a separate one, not a limit.
\item[(d)] With each of the \(2j\) strands of a puncture in the register state of its frame, the count of P follows exactly (maximally mixed
register, all \(2j+1\) states, \(0.274067\)), and R gives one.  The register of M gives a spin-dependent number (tending to \(2.40\) Shannon,
\(1.865\) participation), not the constant two.  The rule of M is equivalent to pure strands with unknown port (one bit per puncture).
\end{enumerate}
Reading: the counts of R and P follow from the register; the count of M does not, and OP1 stays open at exactly this point.

\subsection*{v4.10.0-r28 frames}
\label{subsec:v4100-r28-frames}

\begin{itemize}[leftmargin=2em]
\item M: weight two per spin, \(\gamma_M=0.237533\); used in H102 and in \(\gamma_L\).  Not derived from the register of M.
\item T: \(j=\tfrac12\) only, \(\gamma_T=0.127384\) (ABK).  Branch shifted by \(i\pi/2\) in rapidity.
\item P: all \(2j+1\), \(\gamma_P=0.274067\); a setting, supported by the maximally mixed register.  Common endpoint of both branches.
\end{itemize}

\subsection*{v4.10.0-r28 register entries}
\label{subsec:v4100-r28-register}

\begin{center}
\small
\begin{tabular}{@{}>{\raggedright\arraybackslash}p{0.34\textwidth}>{\raggedright\arraybackslash}p{0.58\textwidth}@{}}
\toprule
statement & status after v4.10.0-r28\\
\midrule
residual bend of the lepton phase (OP3) & open; carrier variants H106--H111 refuted; chain ends at the counting rule of M\\
H102 (\(m_\tau=1777.134\) MeV) & open prediction, \(+2.3\sigma\); joint test of core reading, U(1) counting and spectrum \(\sqrt{j(j+1)}\)\\
counting rule of frame M (OP1) & open; native U(1) axis found, weight two not derived\\
(S1), (A) & (A) needs a lock between two SU(2); (S1\('\)) gives the weight two but breaks (A) in its r24 form\\
count of P & all \(2j+1\) from the register (under (S1)); setting, not compared with data\\
\bottomrule
\end{tabular}
\end{center}

\subsection*{v4.10.0-r28 notes}
\label{subsec:v4100-r28-notes}

Pre-registrations and results H106--H111 with scripts k228--k233: \texttt{notes/\allowbreak research/\allowbreak predictions/}.  Explorations:
\texttt{op3\_\allowbreak renormalisation}, \texttt{op3\_\allowbreak frame\_\allowbreak typed\_\allowbreak halves},
\texttt{op3\_\allowbreak frame\_\allowbreak algebras\_\allowbreak observer}, \texttt{op3\_\allowbreak dixon\_\allowbreak times\_\allowbreak sedenions},
\texttt{op3\_\allowbreak native\_\allowbreak kappa}, \texttt{op3\_\allowbreak chain\_\allowbreak diagnosis},
\texttt{op3\_\allowbreak correction\_\allowbreak term}, \texttt{op3\_\allowbreak topology\_\allowbreak core},
\texttt{op1\_\allowbreak counting\_\allowbreak rule\_\allowbreak data}, \texttt{op1\_\allowbreak u1\_\allowbreak check},
\texttt{op1\_\allowbreak s1\_\allowbreak check}, \texttt{op1\_\allowbreak r\_\allowbreak readout\_\allowbreak precheck},
\texttt{op1\_\allowbreak hc\_\allowbreak blur} (all dated 2026-11-16) and \texttt{op1\_\allowbreak entropy\_\allowbreak count} (2026-11-17), each under
\texttt{notes/\allowbreak research/}.  Except H106--H111, all are explorations and not blind: targets and weights were known.

\subsection*{v4.10.0-r28 ledger and guard update}
\label{subsec:v4100-r28-ledger-guard}

\[
R_{\rm chk}^{v4100\text{-}r28}=(1,1,0)
\]
for (status recorded; computations checked in the cited scripts; the counting rule of frame M and the size of the core not derived).

The layer spends no guard:
\[
R_{\rm glob}^{(v4100\text{-}r28,+)}=(0,1,1,1,1),
\qquad
\Gamma_{\rm guard}^{v4100\text{-}r28}=(1,1,1,1,1),
\qquad
\beta_4=1/256\ \text{unspent}.
\]
The layer does not imply exact finite \(C5\) zero closure, exact all-readout terminality, \(C6/T2\), a global photon-clock/RH theorem,
or global all-readout closure.  OP1--OP5, C5 and C6/T2 remain open.
